Repository instruction files such as AGENTS.md guide AI coding agents, and a nested file may legitimately refine or
replace guidance from its ancestors, so an apparent disagreement is not necessarily a defect. CIPHER defines a
task-conditioned classification of instruction pairs under a declared directory-precedence profile into five labels:
unresolved conflict, valid override, disjoint applicability, compatible, and insufficient evidence. From a public
configuration dataset we built a traceable pilot set of \NOriginals{} real instruction pairs from \NEvalReposWord{}
repositories at pinned commits, each with a context-changing variant and an order control, and a blinded
annotation protocol. Human reference labels have not yet been collected, so this short paper reports a
preliminary analysis, frozen by hash before predictions were opened, that needs none: accuracy on the \NPDMain{}
originals and variants whose label the profile fixes from file paths, and label-free measures elsewhere. With one local
open-weight model, a direct prompt given the full evidence and a structured pipeline that computes applicability and
precedence explicitly were both correct on all \NPDMain{} profile-determined cases, whereas a keyword heuristic flagged
\BZeroPDFalseConf{} of \NPD{} such cases as conflicts. No model judged any original or variant to be an unresolved
conflict; the methods disagreed mainly on condition-based applicability, and swapping only the order of the two
instructions changed 4 of \NControls{} predictions for each model arm. The structured pipeline used more output tokens
than the direct prompt (\BTwoMeanOut{} versus \BOneMeanOut{} per case) without a measurable gain. Human labels for the
semantic cases are the next step.
